\subsection{Ablation and Cross-Cohort Analysis}
\label{sec:ablation}

\paragraph{Algorithm extensions (quick ablation).}
Table~\ref{tab:algo-ablation} compares fixed-threshold B2 against dynamic $\tau_t$ (B2d), multi-frame EMA (B2m), and utility routing (B2u) on 60 ShezhenV3 test frames (seed~0, physical resample).
B2m and B2u raise M2 from 0.60 to 0.65 at the cost of extra edge compute (B2m) or stricter upload utility (B2u); B2d is slightly more conservative under simulated packet loss.
Full 3-seed sweeps and network perturbation curves are scheduled for the supplementary track.

\begin{table}[t]
\centering
\caption{Quick ablation of routing extensions (seed~0, $n{=}60$, physical resample).}
\label{tab:algo-ablation}
\small
\begin{tabular}{lrrr}
\hline
Baseline & M1 p50 (ms) & M2 & M3 \\
\hline
B2 (fixed $\tau$) & 211.3 & 0.600 & 0.417 \\
B2d (dynamic $\tau_t$) & 209.0 & 0.550 & 0.467 \\
B2m (EMA fusion) & 280.4 & \textbf{0.650} & 0.417 \\
B2u (utility $U$) & 205.4 & \textbf{0.650} & 0.367 \\
\hline
\end{tabular}
\end{table}

\paragraph{Threshold $\tau$ (M4).}
We sweep five values of $\tau$ for baseline B2 (500 frames, seed~0) using real Edge-IQA scores from the ShezhenV3 test plan.
Figure~\ref{fig:ablation-tau} shows valid-frame rate peaks near $\tau^\ast{=}0.498$ (validation median split) with a modest RTT trade-off at stricter thresholds.

\begin{figure}[t]
  \centering
  \includegraphics[width=0.95\linewidth]{fig7_ablation_tau.pdf}
  \caption{Ablation over routing threshold $\tau$ (M4): valid rate and RTT p50 for B2 on TCM test frames.}
  \label{fig:ablation-tau}
\end{figure}

\paragraph{Cross-cohort correlation (M5).}
On the full ShezhenV3 validation split ($n{=}1{,}144$ clear/blur pairs), Spearman correlation between $Q_{\mathrm{img}}$ and binary labels is $\rho{=}0.363$ (\texttt{cross\_tcm\_fd.csv}).
This is lower than the synthetic phase-2 cohort ($\rho{\approx}0.75$) but confirms monotonic separation on average; absolute $\tau$ calibration uses cohort medians rather than assuming disjoint score ranges.

\paragraph{Auxiliary Franka track (M6).}
Fifty closed-loop trials on the execution stack produce Fig.~S1; local loop latency does not include simulated cloud upload and must not be compared numerically to Table~\ref{tab:main-results}.

\paragraph{Reproducibility.}
Dataset import, $\tau$ calibration, matrix replay, and figure generation are orchestrated by \texttt{scripts/paper1\_run\_tcm\_full.sh} with paths configured in \texttt{experiments/configs/tcm\_paths.yaml}.
